%; whizzy chapter -dvi
% -initex iniptex -latex platex -format platex -bibtex jbibtex -fmt fmt
% 以上 whizzytex を使用する場合の設定。

%     Tokyo Debian Meeting resources
%     Copyright (C) 2012 Junichi Uekawa
%     Copyright (C) 2011, 2015 Nobuhiro Iwamatsu

%     This program is free software; you can redistribute it and/or modify
%     it under the terms of the GNU General Public License as published by
%     the Free Software Foundation; either version 2 of the License, or
%     (at your option) any later version.

%     This program is distributed in the hope that it will be useful,
%     but WITHOUT ANY WARRANTY; without even the implied warranty of
%     MERCHANTABILITY or FITNESS FOR A PARTICULAR PURPOSE.  See the
%     GNU General Public License for more details.

%     You should have received a copy of the GNU General Public License
%     along with this program; if not, write to the Free Software
%     Foundation, Inc., 51 Franklin St, Fifth Floor, Boston, MA  02110-1301 USA

%  preview (shell-command (concat "evince " (replace-regexp-in-string "tex$" "pdf"(buffer-file-name)) "&"))

%%ここからヘッダ開始。

\documentclass[mingoth,a4paper]{jsarticle}
\usepackage{monthlyreport}
% 日付を定義する、毎月変わります。
\newcommand{\debmtgyear}{2019}
\newcommand{\debmtgmonth}{5}
\newcommand{\debmtgdate}{18}
% started from zero:
% (let ((year 2013) (month 7)) (+ (* (- year 2005) 12) month -1))
\newcommand{\debmtgnumber}{174}

% Needed to import pandoc-generated LaTeX documents.
% See https://stackoverflow.com/questions/40438037/tightlist-error-using-pandoc-with-markdown
\providecommand{\tightlist}{%
  \setlength{\itemsep}{0pt}\setlength{\parskip}{0pt}}

\begin{document}

\begin{titlepage}
\thispagestyle{empty}
% タイトルページ:編集必要な部分は最初のマクロに飛ばすこと

\vspace*{-2cm}
第\debmtgnumber{}回 東京エリア Debian 勉強会資料\\
\hspace*{-2cm}
\includegraphics{image2012-natsu/dotdeb.pdf}\\
\hfill{}\debmtgyear{}年\debmtgmonth{}月\debmtgdate{}日

% ここはアップデートすること
% 全角文字にしないとフォントのサイズが合わないので注意
\rotatebox{10}{\fontsize{30}{30} {\gt　/usr Merge 特集}}\\

\vspace*{-2cm}
\hfill{}\includegraphics[height=6cm]{image200502/openlogo-nd.eps}
\end{titlepage}

\newpage

\begin{minipage}[b]{0.2\hsize}
 \definecolor{titleback}{gray}{0.9}
 \colorbox{titleback}{\rotatebox{90}{\fontsize{80}{80} {\gt デビアン勉強会} }}
\end{minipage}
\begin{minipage}[b]{0.8\hsize}
\hrule
\vspace{2mm}
\hrule
\begin{multicols}{2}
\tableofcontents
\end{multicols}
\vspace{2mm}
\hrule
\end{minipage}

\dancersection{最近のDebian関連のミーティング報告}{杉本　典充}

\subsection{第173回東京エリアDebian勉強会}

2019年4月20日(土)に第173回東京エリアDebian勉強会を開催しました。会場は東銀座にある朝日ネットさんをお借りして行いました。参加者は5名でした。

セミナー発表は、野首さんの「grml-debootstrapを用いたUSB起動メモリの作成」でした。USBメモリにDebianをインストールして起動する使い方の便利さやUSBメモリの作り方や作るときに考慮するとよい点を説明しました。

その後Hack Timeを行い、各自が用意した課題に取り組みました。


\dancersection{事前課題}{dictoss}

今回の事前課題は以下です。

\begin{enumerate}
%\item xxx
\item Hack Timeは何をしますか (How will you work on Hack Time ?)
\end{enumerate}

%この課題に対して提出いただいた内容は以下です。

\begin{multicols}{2}
{\small
\begin{prework}{ yy\_y\_ja\_jp }
  \begin{enumerate}
  \item $B!J2sEz$J$7!K(B
  \end{enumerate}
\end{prework}

\begin{prework}{ khibino }
  \begin{enumerate}
  \item $B!J2sEz$J$7!K(B
  \end{enumerate}
\end{prework}

\begin{prework}{ NOKUBI Takatsugu (knok) }
  \begin{enumerate}
  \item $B%a%$%s$r(Bbuster$B4D6-$K$7$?$$(B
  \end{enumerate}
\end{prework}

\begin{prework}{ lion (2015fuj) }
  \begin{enumerate}
  \item $B!J2sEz$J$7!K(B
  \end{enumerate}
\end{prework}

\begin{prework}{ Kouhei Maeda (mkouhei) }
  \begin{enumerate}
  \item $B!J2sEz$J$7!K(B
  \end{enumerate}
\end{prework}

\begin{prework}{ dictoss }
  \begin{enumerate}
  \item Debian$B$N3+H/<T8~$1%I%-%e%a%s%H$rFI$`(B
  \end{enumerate}
\end{prework}

}
\end{multicols}

%\dancersection{Debian Trivia Quiz}{username}
%
%Debianの昨今の話題についてのQuizです。
%
%今回の出題範囲は\url{debian-devel-announce@lists.debian.org} や \url{debian-news@lists.debian.org}などに投稿された内容からです。
%
%\begin{multicols}{2}
%%; whizzy-master ../debianmeetingresume201211.tex
% $B0J>e$N@_Dj$r$7$F$$$k$?$a!"$3$N%U%!%$%k$G(B M-x whizzytex $B$9$k$H!"(Bwhizzytex$B$,MxMQ$G$-$^$9!#(B
%

\santaku
{DebConf13 $B$N3+:ECO$H3+:EF|$O!)(B}
{$BF|K\!"El5~ET(B 6$B7n(B20$BF|(B}
{$B%K%+%i%0%"(B $B%^%J%0%"(B 7$B7n(B8-14$BF|(B}
{$B%9%$%9!"%t%)!<%^%k%-%e(B 8$B7n(B11-18$BF|(B}
{3}
{$B%K%+%i%0%"$O(BDebConf12$B$N3+:ECO$G$9!#(B
DebConf13$B$O%9%$%9$N%-%c%s%WCO$G3+:E$G$9!#(B
6/20$B$O3'$5$sM=Dj$r6u$1$F$*$-$^$7$g$&!#(B}

\santaku
{$B@$3&$N(BWeb$B%5!<%P$G:G$b?M5$$N$"$k(BLinux $B%G%#%9%H%j%S%e!<%7%g%s(B(W3Techs$BD4$Y(B)$B$O!)(B}
{CentOS}
{Debian}
{Ubuntu}
{B}
{\url{http://w3techs.com/technologies/history_details/os-linux}$B$K7k2L$N%0%i%U$,$"$j$^$9!#(B
$B8=:_(B Linux $B$r;HMQ$7$F$$$k(B web $B%5!<%P$N(B 32.9\% $B$,(B Debian $B$rMxMQ$7$F$*$j!"$=$N3d9g$O8=:_$bA}2C$rB3$1$F$$$k$=$&$G$9!#(B}

\santaku
{Debian $B%+!<%M%k%A!<%`$N%a%s%P!<$G$"$j!"(Bkernel.org $B$N(B 3.2.y $B0BDjHG7ONs$N%a%s%F%J$G$b$"$k(B Ben Hutchings $B$5$s$,<!4|(B Debian $B0BDjHG$H0l=o$K=P2Y$5$l$k(B Linux $B%+!<%M%k$K(B (3.2 $B7ONs$N(B mainline $B$K$OL5$$(B) $BDI2C5!G=$,Ek:\$5$l$kM=Dj$G$"$k$H=R$Y$F$$$^$9!#(B
$BB?$/$NDI2CE@$NCf$K4^$^$l$J$$$b$N$O2?!)(B}
{PREEMPT\_RT}
{Hyper-V guest drivers$B$N6/2=(B}
{ARM64/AArch64$B%"!<%-%F%/%A%c%5%]!<%H(B}
{C}
{Hyper-V guest drivers$B$O(Bmainline kernel$B$G(B3.2$B$K$b4^$^$l$F$$$^$9$,!"$h$j2~A1$5$l$?(B3.4$B$+$i$N=$@5$,F3F~$5$l$^$9!#(B
PREEMPT\_RT$B$O%O!<%I%j%"%k%?%$%`$r<B8=$9$k$?$a$N(BPatch$B!"(B
linux-image-rt-amd64 , linux-image-rt-686-pae $B$N(Bmetapackage$B$G;HMQ$G$-$^$9!#(B
$B?7$7$$(BARM 64$B%S%C%H%"!<%-%F%/%A%c%5%]!<%H$O(Bmainline kernel 3.7$B$+$i(B}

\santaku
{Wookey$B$5$s$,%"%J%&%s%9$7$?(Balpha$BHG$N(BDebian port arm64 image$B$O!)(B}
{Debian/Ubuntu port image}
{Debian/KFreeBSD port image}
{Debian/GnuHurd port image}
{A}
{self-bootstrapp(non x86)$BBP1~$H$N$3$H$G$9!#(B\url{http://wiki.debian.org/Arm64Port}$B$G%9%F!<%?%9$,3NG'$G$-$^$9!#(B}

\santaku
{700,000$BHVL\$N%P%0$,Js9p$5$l$?F|$rEv$F$k(B700000thBugContest$B$N7k2L$,=P$^$7$?!#$=$NM=A[F|$HJs9pF|$O!)(B}
{2012/12/12$B$rM=A[$7$?(BDavidPrevot}
{$BM=A[F|(B:2013/02/04$B!"Js9pF|(B:2013/02/14}
{$BM=A[F|(B:2013/02/07$B!"Js9pF|(B:2013/02/14}
{$BM=A[F|(B:2013/02/14$B!"Js9pF|(B:2013/02/07}
{C}
{$B:G$b6a$$(B2013/02/14$B$rM=A[$7$?(BChristian Perrier$B$5$s$,Ev$F$^$7$?!#7k2L$O(B\url{http://wiki.debian.org/700000thBugContest}$B$G8x3+$5$l$F$$$^$9!#(B
$B$^$?!"(B800,000/1,000,000$BHVL\$N%P%0$,Js9p$5$l$kF|$rEv$F$k%3%s%F%9%H(B\url{http://wiki.debian.org/800000thBugContest}$B$b3+:E$5$l$F$$$^$9!#(B}

\santaku
{master.debian.org$B$,?7$7$$5!3#$K0\9T$5$l$^$7$?!#$3$l$O2?$N%5!<%P$G$7$g$&$+(B $B!)(B}
{@debian.org$B$N%a!<%k%5!<%P(B}
{$B%Q%C%1!<%8$N%^%9%?!<%5!<%P(B}
{$B%Q%C%1!<%8$N%9%]%s%5!<(B(mentor)$B$rC5$9%5!<%P(B}
{A}
{$B8E$$%5!<%P$O%G%#%9%/>c32Ey$,$"$C$?$N$G!"<wL?$HH=CG$5$l!"%G!<%?$,B;<:$9$kA0$K?7$7$$%5!<%P$K0\9T$5$l$^$7$?!#(Bftp-master.debian.org$B$O(BDebian$B$N(B official package $B%j%]%8%H%j$G$9!#%Q%C%1!<%8$N%9%]%s%5!<(B(mentor)$B$rC5$9$N$O(Bmentors.debian.net$B!#(B }

\santaku
{pbuilder$B$K(Bclang support$B$,DI2C$5$l$^$7$?!#C/$,=q$$$?%Q%C%A$G$7$g$&$+!)(B}
{Sylvestre Ledru}
{Junichi Uekawa}
{Hideki Yamane}
{C}
{Debian$B$N(BClang$B%5%]!<%H$OCe!9$H?J$s$G$$$^$9!#(B}

\santaku
{DPN - 2013$BG/(B3$B7n(B4$BF|9f$K<h$j>e$2$i$l$?F|K\$N%$%Y%s%H$O(B}
{Open Source Conference 2013 Tokyo/Spring}
{Open Source Conference 2013 Hamamatu}
{Open Source Conference 2013 Tokushima}
{A}
{\url{http://henrich-on-debian.blogspot.jp/2013/02/open-source-conference-2013-tokyospring.html} $B>\:Y$O8e$[$I!#(B}


%\end{multicols}


% % (query-replace-regexp "<.*?>" "")
% % (query-replace-regexp "^[    ]\+" "")

%-------------------------------------------------------------------------------
\dancersection{/usr Mergeについて}{yy\_y\_ja\_jp}
%-------------------------------------------------------------------------------

Debian では Buster から ``/usr Merge'' がデフォルトになりそうです。概要と状況を整理します。

\subsection{概要}

``/usr Merge'' とは、\verb|/bin/|, \verb|/sbin/|, \verb|/lib/| などの中にあるもの（バイナリ等）を \verb|/usr/| 内に移す、つまり \verb|/usr/| 内にバイナリ等をすべて統合することです。
例えば \verb|/bin/ls| を \verb|/usr/bin/ls| に移します。また、今まで通り \verb|/bin/ls| でも使えるように互換性のため \verb|/bin| は \verb|/usr/bin| へのシンボリックリンクにします。\verb|/sbin| なども同様です。

なお、このような構成にしたい理由はsystemdのwiki\footnote{\url{https://www.freedesktop.org/wiki/Software/systemd/TheCaseForTheUsrMerge/}}などにまとめられているようです
\footnote{Debian wiki \url{https://wiki.debian.org/UsrMerge}}。

Debianを新規インストールするときには、Debianインストーラが内部で実行するdebootstrapがこの構成を行います。既存のDebian環境にはusrmergeパッケージをインストールすればこの構成にしてくれます\footnote{Busterの次のバージョンBullseyeではインストール推奨となるかもしれません。\url{https://bugs.debian.org/841666}}。なお、以前の構成へ元に戻すことは困難です。

Stretch インストーラで作った Debian 環境の \verb|/| 直下は次のようなディレクトリ構成でした。
\begin{commandline}
$ ls -l /
合計 76
drwxr-xr-x  2 root root  4096  5月 14 00:42 bin
drwxr-xr-x  3 root root  4096  5月 14 00:43 boot
drwxr-xr-x 17 root root  2980  5月 14 00:56 dev
drwxr-xr-x 76 root root  4096  5月 14 00:56 etc
drwxr-xr-x  3 root root  4096  5月 14 00:43 home
lrwxrwxrwx  1 root root    29  5月 14 00:38 initrd.img -> boot/initrd.img-4.9.0-9-amd64
lrwxrwxrwx  1 root root    29  5月 14 00:38 initrd.img.old -> boot/initrd.img-4.9.0-9-amd64
drwxr-xr-x 15 root root  4096  5月 14 00:42 lib
drwxr-xr-x  2 root root  4096  5月 14 00:36 lib64
drwx------  2 root root 16384  5月 14 00:35 lost+found
drwxr-xr-x  3 root root  4096  5月 14 00:35 media
drwxr-xr-x  2 root root  4096  5月 14 00:35 mnt
drwxr-xr-x  2 root root  4096  5月 14 00:35 opt
dr-xr-xr-x 80 root root     0  5月 14  2019 proc
drwx------  2 root root  4096  5月 14 00:35 root
drwxr-xr-x 14 root root   460  5月 14 00:57 run
drwxr-xr-x  2 root root  4096  5月 14 00:56 sbin
drwxr-xr-x  2 root root  4096  5月 14 00:35 srv
dr-xr-xr-x 13 root root     0  5月 14 00:57 sys
drwxrwxrwt  8 root root  4096  5月 14 00:56 tmp
drwxr-xr-x 10 root root  4096  5月 14 00:35 usr
drwxr-xr-x 11 root root  4096  5月 14 00:35 var
lrwxrwxrwx  1 root root    26  5月 14 00:38 vmlinuz -> boot/vmlinuz-4.9.0-9-amd64
lrwxrwxrwx  1 root root    26  5月 14 00:38 vmlinuz.old -> boot/vmlinuz-4.9.0-9-amd64
$ 
\end{commandline}

一方、Buster インストーラで作った Debian 環境では次のようになります。なお、現時点ではインストーラに ``/usr Merge'' しないようにする設定はなさそうです\footnote{\url{https://bugs.debian.org/923091}}。
\begin{commandline}
$ ls -l /
合計 60
lrwxrwxrwx  1 root root     7  5月 14 00:36 bin -> usr/bin
drwxr-xr-x  3 root root  4096  5月 14 00:43 boot
drwxr-xr-x 17 root root  3160  5月 14 00:56 dev
drwxr-xr-x 67 root root  4096  5月 14 00:56 etc
drwxr-xr-x  3 root root  4096  5月 14 00:43 home
lrwxrwxrwx  1 root root    30  5月 14 00:38 initrd.img -> boot/initrd.img-4.19.0-4-amd64
lrwxrwxrwx  1 root root    30  5月 14 00:38 initrd.img.old -> boot/initrd.img-4.19.0-4-amd64
lrwxrwxrwx  1 root root     7  5月 14 00:36 lib -> usr/lib
lrwxrwxrwx  1 root root     9  5月 14 00:36 lib32 -> usr/lib32
lrwxrwxrwx  1 root root     9  5月 14 00:36 lib64 -> usr/lib64
lrwxrwxrwx  1 root root    10  5月 14 00:36 libx32 -> usr/libx32
drwx------  2 root root 16384  5月 14 00:36 lost+found
drwxr-xr-x  3 root root  4096  5月 14 00:36 media
drwxr-xr-x  2 root root  4096  5月 14 00:36 mnt
drwxr-xr-x  2 root root  4096  5月 14 00:36 opt
dr-xr-xr-x 79 root root     0  5月 14  2019 proc
drwx------  2 root root  4096  5月 14 00:36 root
drwxr-xr-x 15 root root   460  5月 14 00:58 run
lrwxrwxrwx  1 root root     8  5月 14 00:36 sbin -> usr/sbin
drwxr-xr-x  2 root root  4096  5月 14 00:36 srv
dr-xr-xr-x 13 root root     0  5月 14 00:56 sys
drwxrwxrwt  8 root root  4096  5月 14 00:56 tmp
drwxr-xr-x 13 root root  4096  5月 14 00:36 usr
drwxr-xr-x 11 root root  4096  5月 14 00:36 var
lrwxrwxrwx  1 root root    27  5月 14 00:38 vmlinuz -> boot/vmlinuz-4.19.0-4-amd64
lrwxrwxrwx  1 root root    27  5月 14 00:38 vmlinuz.old -> boot/vmlinuz-4.19.0-4-amd64
$ 
\end{commandline}

\subsection{経緯}

``/usr Merge'' は以前Fedoraで特に行われていたようです\footnote{\url{https://fedoraproject.org/wiki/Features/UsrMove}}。なおFedoraは必須化したようです。

Debianでは当初Stretchでデフォルトにする予定でしたが、問題の解決が間に合いそうになく延期されました\footnote{\url{https://lists.debian.org/debian-devel-announce/2017/01/msg00004.html}}。

その後Busterに向けてdebootstrap 1.0.102で再度デフォルトにすることにしました\footnote{\url{https://bugs.debian.org/839046}}が、問題があるためデフォルトでなくすべきではという議論があり、Debian技術委員会(tech-ctte)にも持ち込まれました\footnote{\url{https://bugs.debian.org/914897}}。

%
Debian技術委員会は、debootstrapのメンテナーが決めたことは覆さない、
%つまり強権発動はしない
と決定しました\footnote{\url{https://lists.debian.org/debian-devel-announce/2019/03/msg00001.html}}。
同時に、Debianにおける``/usr Merge''の現状と今後の望ましい姿を示しました。

\subsection{現状}
Debian技術委員会の決定により、debootstrapメンテナー次第で``/usr Merge''がデフォルトになるかどうかは決まります。現時点で testing (Buster) 以降のバージョン向けにdebootstrapを実行すると次のようになっており、メンテナーが変更しなければこのままBusterリリースとなります。

\begin{itemize}
 \item 通常の環境では``/usr Merge''がデフォルト有効（無効にするには \verb|--no-merged-usr|オプションを付ける）
 \item パッケージビルド用の環境（\verb|--variant=buildd| オプション）ではデフォルト無効
\end{itemize}

なお、現状では``/usr Merge''された環境でビルドされたバイナリパッケージを``/usr Merge''されていない環境に持っていくと動作しないものがあり、それらのパッケージに随時バグ報告してくれている方がいるようです\footnote{\url{https://bugs.debian.org/cgi-bin/pkgreport.cgi?tag=usrmerge;users=md@linux.it}}
。よくある理由としては、ビルド中にバイナリの場所を \verb|/bin/| の中よりも先に \verb|/usr/bin/| を探しに行き、``/usr Merge''された環境では \verb|/usr/bin/| のパスをコンパイル結果に埋め込むからです。
例えば quilt パッケージの少し古い修正\footnote{\url{https://bugs.debian.org/913226}}以前のバージョン 0.65-2 を``/usr Merge''された環境でビルドしてみます。\verb|dh_auto_configure| 経由で実行された \verb|./configure| が \verb|/bin/bash| ではなく \verb|/usr/bin/bash|、\verb|/bin/cp| ではなく \verb|/usr/bin/cp|、などを検出してビルドしています。
この場合はビルドされたパッケージ内のシェルスクリプトの先頭が \verb|#! /usr/bin/bash| となったため lintian がたまたま別件のポリシー違反として検出してくれています。
\begin{commandline}
$ dget http://snapshot.debian.org/archive/debian/20180809T030926Z/pool/main/q/quilt/quilt_0.65-2.dsc
dget: retrieving http://snapshot.debian.org/archive/debian/20180809T030926Z/pool/main/q/quilt/quilt_0.65-2.dsc
(snip)
quilt_0.65-2.dsc:
      Good signature found
   validating quilt_0.65.orig.tar.gz
   validating quilt_0.65-2.debian.tar.xz
All files validated successfully.
dpkg-source: info: extracting quilt in quilt-0.65
(snip)
dpkg-source: info: applying fix-mail-threading
$ cd quilt-0.65/
$ debuild -us -uc
 dpkg-buildpackage -us -uc -ui
dpkg-buildpackage: info: source package quilt
dpkg-buildpackage: info: source version 0.65-2
(snip)
dh_auto_configure -- --with-docdir=/usr/share/doc/quilt --with-sendmail=/usr/sbin/sendmail --with-awk=/usr/bin/awk
	./configure --build=x86_64-linux-gnu --prefix=/usr --includedir=\${prefix}/include --mandir=\${prefix}/share/man --infodir=\${prefix}/share/info --sysconfdir=/etc --localstatedir=/var --disable-silent-rules --libdir=\${prefix}/lib/x86_64-linux-gnu --libexecdir=\${prefix}/lib/x86_64-linux-gnu --runstatedir=/run --disable-maintainer-mode --disable-dependency-tracking --with-docdir=/usr/share/doc/quilt --with-sendmail=/usr/sbin/sendmail --with-awk=/usr/bin/awk
configure: WARNING: unrecognized options: --disable-silent-rules, --disable-maintainer-mode, --disable-dependency-tracking, --with-docdir
checking for a BSD-compatible install... /usr/bin/install -c
checking whether #! works in shell scripts... yes
checking for bash... /usr/bin/bash
checking whether /usr/bin/bash quoting works... yes
checking for gcp... no
checking for cp... /usr/bin/cp
checking for gdate... no
checking for date... /usr/bin/date
(snip)
dpkg-deb: building package 'quilt' in '../quilt_0.65-2_all.deb'.
dpkg-deb: building package 'quilt-el' in '../quilt-el_0.65-2_all.deb'.
 dpkg-genbuildinfo
 dpkg-genchanges  >../quilt_0.65-2_amd64.changes
dpkg-genchanges: info: not including original source code in upload
 dpkg-source --after-build .
dpkg-source: warning: Testsuite field contains value autopkgtest, but no tests control file debian/tests/control
dpkg-buildpackage: info: binary and diff upload (original source NOT included)
Now running lintian quilt_0.65-2_amd64.changes ...
W: quilt source: orig-tarball-missing-upstream-signature quilt_0.65.orig.tar.gz
E: quilt: missing-depends-on-sensible-utils usr/share/quilt/edit
E: quilt: missing-depends-on-sensible-utils usr/share/quilt/header
E: quilt: missing-depends-on-sensible-utils usr/share/quilt/mail
E: quilt: wrong-path-for-interpreter usr/bin/quilt (#!/usr/bin/bash != /bin/bash)
E: quilt: wrong-path-for-interpreter usr/share/quilt/add (#!/usr/bin/bash != /bin/bash)
(snip)
Finished running lintian.
$ 
\end{commandline}

問題を回避するため、少なくともDebian公式のパッケージビルドには当面``/usr Merge''されていない環境を使うこととなりました。また、Debian公式レポジトリにアップロードされたパッケージは、``/usr Merge''されていない環境でビルドした結果と``/usr Merge''された環境でビルドした結果を比較してこのような差分が発生していないかReproducible BuildsプロジェクトのCI\footnote{Continuous Integration、継続的インテグレーション}環境で検証も行われており\footnote{\url{https://tests.reproducible-builds.org/debian/issues/unstable/paths_vary_due_to_usrmerge_issue.html}}、これも元にバグ報告・修正が行われています。

手元で動くpbuilder, cowbuilderやsbuildなどクリーンルームでのパッケージビルドツールも同様に``/usr Merge''されていない環境が作られてその中でビルドしてくれます。新たにcowbuilderでビルド環境を作ると次のように構成されます。

\begin{commandline}
# cowbuilder --create
I: Invoking pbuilder
I: forking: pbuilder create --buildplace /var/cache/pbuilder/base.cow --mirror http://ftp.jp.debian.org/debian --distribution sid --no-targz --extrapackages cowdancer
W: /root/.pbuilderrc does not exist
I: Running in no-targz mode
（略）
Processing triggers for libc-bin (2.28-10) ...
I: Copying back the cached apt archive contents
I: unmounting dev/ptmx filesystem
I: unmounting dev/pts filesystem
I: unmounting dev/shm filesystem
I: unmounting proc filesystem
I: unmounting sys filesystem
# 
\end{commandline}
\begin{commandline}
# ls -l /var/cache/pbuilder/base.cow/
合計 80
drwxr-xr-x  2 root root 4096  5月 14 23:50 bin
drwxr-xr-x  2 root root 4096  5月 14 05:25 boot
drwxr-xr-x  2 root root 4096  5月 14 23:50 build
drwxr-xr-x  4 root root 4096  5月 14 23:50 dev
drwxr-xr-x 30 root root 4096  5月 14 23:50 etc
drwxr-xr-x  2 root root 4096  5月 14 05:25 home
drwxr-xr-x  7 root root 4096  5月 14 23:50 lib
drwxr-xr-x  2 root root 4096  5月 14 23:50 lib64
drwxr-xr-x  2 root root 4096  5月 14 23:50 media
drwxr-xr-x  2 root root 4096  5月 14 23:50 mnt
drwxr-xr-x  2 root root 4096  5月 14 23:50 opt
drwxr-xr-x  2 root root 4096  5月 14 05:25 proc
drwx------  2 root root 4096  5月 14 23:50 root
drwxr-xr-x  4 root root 4096  5月 14 23:50 run
drwxr-xr-x  2 root root 4096  5月 14 23:50 sbin
drwxr-xr-x  2 root root 4096  5月 14 23:50 srv
drwxr-xr-x  2 root root 4096  5月 14 05:25 sys
drwxrwxrwt  2 root root 4096  5月 14 23:50 tmp
drwxr-xr-x 10 root root 4096  5月 14 23:50 usr
drwxr-xr-x 11 root root 4096  5月 14 23:50 var
#
\end{commandline}
パッケージを作るときは安全のためこれらのツールでビルドするようにしましょう。

なお、多くの場合 \verb|PATH| 環境変数に従ってバイナリの場所を探しているので \verb|/bin/| よりも \verb|/usr/bin/| が優先されています。なのでもともと``/usr Merge''されていない環境であってもクリーンルームでビルドしていないと次のようなことも起こります。\verb|/usr/local/bin/bash| が存在する環境で \verb|dpkg-buildpackage| コマンドを実行すると、\verb|PATH| 環境変数では \verb|/bin|（や \verb|/usr/bin|）よりも \verb|/usr/local/bin| が優先されるので \verb|/bin/bash| ではなく \verb|/usr/local/bin/bash| が使われてしまいます。なおこのように \verb|/usr/local/bin/| の中のものを使うのは明らかにポリシー違反なので一部ツールでは対策されており、\verb|dpkg-buildpackage| コマンドの代わりにdevscriptsパッケージにある \verb|debuild| コマンドを使えば回避できます。これは \verb|PATH| 環境変数に \verb|/usr/local/bin| を除いたものを設定した上でビルドしてくれます。ですがまずは先ほどのクリーンルームなビルドツールを使うのが良いでしょう。

また、``/usr Merge''された環境（に限らないですが）の既知の注意点として、\verb|dpkg --search| (\verb|dpkg -S|) が指定したファイルのパッケージを見つけてくれないことがあります\footnote{\url{https://bugs.debian.org/858331}}。
\begin{commandline}
$ which ls
/usr/bin/ls
$ dpkg -S /usr/bin/ls
dpkg-query: パターン /usr/bin/ls に一致するパスが見つかりません
$ dpkg -S /bin/ls
coreutils: /bin/ls
$ 
\end{commandline}
なお、おそらくこの挙動が変更されることはなさそうです。debhelperをはじめ様々なツールも現在の挙動に依存しています。ディレクトリの対応付けなどを認識しながらファイルを検索するオプションは別に用意されるかもしれません。

\subsection{今後}

Debian技術委員会は今後Debianの望ましい姿として、Busterの次のバージョン(Bullseye)では``/usr Merge''されていない環境とされている環境の両方が使えるべきで、またどちらの環境でビルドされたパッケージも両方の環境で動作すべきと示しました。

先ほどの quilt は、その次のバージョン 0.65-3 で``/usr Merge''された環境でビルドされたパッケージが``/usr Merge''されていない環境に持っていっても動作するよう次のように修正されました。debian/rules ファイルの \verb|dh_auto_configure| コマンド（結果的には \verb|./configure|）呼び出しの引数最後に \verb|--with-bash=/bin/bash| を追加し、\verb|PATH| 環境変数に関係なく決め打ちで \verb|/bin/bash| を見に行くようにしました。なお \verb|./configure| は \verb|--with-bash=| を \verb|BASH=| と書くこともできます。今のところこのように差分が見つかった各パッケージを修正中のようです。新しくパッケージングしたいアプリにも少なくとも差分があればこのような対処が必要になります。

さて、現時点でこれをどうシームレスに実現するのかは不明です。少なくともDebian公式レポジトリにアップロードされた後に検出するCI環境は既にありますが、新しいアプリをパッケージングするときに事前に対応できるようにならないと真に実現したことにならないでしょう。そのためには、debhelper等のパッケージング時のビルド補助ツールが頑張り、lintian等のパッケージング後の検査ツールも強化していく感じでしょうか。\verb|PATH| 環境変数で \verb|/bin| よりも \verb|/usr/bin| を優先させることもできるかもしれませんが悪影響が読めません。現時点でこの問題の自衛方法としてはReproducible Buildsの取り組みで使われているツールを手元で動かすくらいでしょう。diffoscopeで2つのファイルをバイナリの中身まで再帰的に差分表示してくれます\footnote{x11-apps で発生している差分の表示例 \url{https://tests.reproducible-builds.org/debian/dbd/buster/amd64/x11-apps_7.7+7.diffoscope.html}。表示に使ったdiffoscopeコマンド引数はこのHTMLページのタイトル(\texttt{<title>})に書かれています。}。ただ2つの環境作って両方でビルドするのも、そしてこの``/usr Merge''で発生している差分を見極めるのも正直面倒なので\footnote{完全なReproducible Buildsは本当は実現すべきですが大変なので、今回発生する差分だけでも対処するのが現実的です。}、まずはクリーンルームでビルドして上げておいて、後でCI環境などに指摘されたら直すくらいが現実的ではないかと思います。

%
% 冊子にするために、4の倍数にする必要がある。
% そのための調整
\dancersection{メモ}{}
\mbox{}\newpage
\mbox{}\newpage
\mbox{}\newpage

\vspace*{15cm}
\hrule
\vspace{2mm}
\includegraphics[width=2cm]{image200502/openlogo-nd.eps}
\noindent \Large \bf Debian 勉強会資料\\
\noindent \normalfont \debmtgyear{}年\debmtgmonth{}月\debmtgdate{}日 \hspace{5mm}  初版第1刷発行\\
\noindent \normalfont 東京エリア Debian 勉強会 （編集・印刷・発行）\\
\hrule

\end{document}
